% !TEX program = xelatex
\documentclass[10pt,a4paper]{article}
\usepackage[utf8]{inputenc}
\usepackage[T1]{fontenc}
\usepackage[french]{babel}
\usepackage[left=0.6cm,right=0.6cm,top=1cm,bottom=1cm]{geometry}
\usepackage{amsmath}
\usepackage{amssymb}
\usepackage{array}
\usepackage{graphicx}
\usepackage{xcolor}
\usepackage{titlesec}
\usepackage{enumitem}
\usepackage{multicol}
\usepackage{pifont}
% Badges « Exercice n » en cursive Playwrite FR Moderne (police du site),
% chargée à 85 % de la taille courante (XeLaTeX ou LuaLaTeX ; gras sous pdfLaTeX).
\makeatletter
\ifdefined\XeTeXrevision\def\ExoFont@ok{}\fi
\ifdefined\directlua\def\ExoFont@ok{}\fi
\ifdefined\ExoFont@ok
\newcommand{\ExoFont}{\font\ExoFont@f="[../polices/PlaywriteFRModerne.ttf]" at \dimexpr\f@size pt*85/100\relax\ExoFont@f}
\else
\newcommand{\ExoFont}{\bfseries}
\fi
\makeatother
% \ExoTexte{...} : texte du badge en cursive, avec la hauteur d'une majuscule de la
% police courante (la cursive est plus haute et ferait grandir chaque badge).
\newcommand{\ExoTexte}[1]{\raisebox{0pt}[\fontcharht\font`E][0pt]{\ExoFont #1}}
\usepackage{frcursive}
% Correctif du bug de crénage de frcursive sous XeLaTeX (repris de CLT2627.sty) :
% sans lui, certaines paires de lettres consécutives dans \cursive explosent
% (une lettre projetée loin à droite), ce qui coupe des mots comme "Représenter".
\makeatletter
\ifdefined\XeTeXcharclass
\XeTeXcharclass `a 201 \XeTeXcharclass `b 201 \XeTeXcharclass `c 201
\XeTeXcharclass `d 201 \XeTeXcharclass `e 201 \XeTeXcharclass `f 201
\XeTeXcharclass `g 201 \XeTeXcharclass `h 201 \XeTeXcharclass `i 201
\XeTeXcharclass `j 201 \XeTeXcharclass `k 201 \XeTeXcharclass `l 201
\XeTeXcharclass `m 201 \XeTeXcharclass `n 201 \XeTeXcharclass `o 201
\XeTeXcharclass `p 201 \XeTeXcharclass `q 201 \XeTeXcharclass `r 201
\XeTeXcharclass `s 201 \XeTeXcharclass `t 201 \XeTeXcharclass `u 201
\XeTeXcharclass `v 201 \XeTeXcharclass `w 201 \XeTeXcharclass `x 201
\XeTeXcharclass `y 201 \XeTeXcharclass `z 201
\XeTeXcharclass `A 201 \XeTeXcharclass `B 201 \XeTeXcharclass `C 201
\XeTeXcharclass `D 201 \XeTeXcharclass `E 201 \XeTeXcharclass `F 201
\XeTeXcharclass `G 201 \XeTeXcharclass `H 201 \XeTeXcharclass `I 201
\XeTeXcharclass `J 201 \XeTeXcharclass `K 201 \XeTeXcharclass `L 201
\XeTeXcharclass `M 201 \XeTeXcharclass `N 201 \XeTeXcharclass `O 201
\XeTeXcharclass `P 201 \XeTeXcharclass `Q 201 \XeTeXcharclass `R 201
\XeTeXcharclass `S 201 \XeTeXcharclass `T 201 \XeTeXcharclass `U 201
\XeTeXcharclass `V 201 \XeTeXcharclass `W 201 \XeTeXcharclass `X 201
\XeTeXcharclass `Y 201 \XeTeXcharclass `Z 201
\XeTeXcharclass `\é 201 \XeTeXcharclass `\è 201 \XeTeXcharclass `\ê 201
\XeTeXcharclass `\ë 201 \XeTeXcharclass `\É 201 \XeTeXcharclass `\È 201
\XeTeXcharclass `\Ê 201 \XeTeXcharclass `\Ë 201 \XeTeXcharclass `\à 201
\XeTeXcharclass `\â 201 \XeTeXcharclass `\ä 201 \XeTeXcharclass `\À 201
\XeTeXcharclass `\Â 201 \XeTeXcharclass `\Ä 201 \XeTeXcharclass `\ô 201
\XeTeXcharclass `\ö 201 \XeTeXcharclass `\Ô 201 \XeTeXcharclass `\Ö 201
\XeTeXcharclass `\î 201 \XeTeXcharclass `\ï 201 \XeTeXcharclass `\Î 201
\XeTeXcharclass `\Ï 201 \XeTeXcharclass `\ù 201 \XeTeXcharclass `\û 201
\XeTeXcharclass `\ü 201 \XeTeXcharclass `\Ù 201 \XeTeXcharclass `\Û 201
\XeTeXcharclass `\Ü 201 \XeTeXcharclass `\ç 201 \XeTeXcharclass `\Ç 201
\XeTeXcharclass `\œ 201 \XeTeXcharclass `\Œ 201 \XeTeXcharclass `\æ 201
\XeTeXcharclass `\Æ 201
\let\cursive@orig\cursive
\renewcommand{\cursive}{%
  \cursive@orig
  \XeTeXinterchartoks 201 201 = {\kern0pt }%
}
\fi
\makeatother
\usepackage{tikz}
\usetikzlibrary{calc,arrows.meta,shapes.misc}
\tikzset{ptcross/.style={cross out,draw,thick,minimum size=3pt,inner sep=0pt}}
\usepackage{tcolorbox}
\tcbuselibrary{skins,breakable}
\usepackage[hidelinks]{hyperref}
\usepackage{fontawesome5}
\definecolor{calcblue}{HTML}{1C7ED6}
\definecolor{diffgreen}{HTML}{40C057}
\definecolor{diffblue}{HTML}{228BE6}
\definecolor{diffred}{HTML}{FA5252}
\definecolor{skyblue}{HTML}{339AF0}
\definecolor{objorange}{HTML}{F76707}
\newcommand{\starfull}[1]{{\color{#1}\ding{72}}}
\newcommand{\starempty}{{\color{gray!45}\ding{73}}}
\newcommand{\Vj}[1]{\overrightarrow{\text{#1}}}
\newcommand{\VBlank}{\rule[0.15ex]{0.55cm}{0.35pt}}
\newcommand{\Pt}[1]{\textsf{\footnotesize #1}}
\renewcommand{\familydefault}{\sfdefault}
\DeclareSymbolFont{sfnumbers}{OT1}{cmss}{m}{n}
\SetSymbolFont{sfnumbers}{bold}{OT1}{cmss}{bx}{n}
\DeclareMathSymbol{0}{\mathalpha}{sfnumbers}{`0}
\DeclareMathSymbol{1}{\mathalpha}{sfnumbers}{`1}
\DeclareMathSymbol{2}{\mathalpha}{sfnumbers}{`2}
\DeclareMathSymbol{3}{\mathalpha}{sfnumbers}{`3}
\DeclareMathSymbol{4}{\mathalpha}{sfnumbers}{`4}
\DeclareMathSymbol{5}{\mathalpha}{sfnumbers}{`5}
\DeclareMathSymbol{6}{\mathalpha}{sfnumbers}{`6}
\DeclareMathSymbol{7}{\mathalpha}{sfnumbers}{`7}
\DeclareMathSymbol{8}{\mathalpha}{sfnumbers}{`8}
\DeclareMathSymbol{9}{\mathalpha}{sfnumbers}{`9}
\setlength{\parindent}{0pt}
\graphicspath{{./}}
\newcommand*{\textcircledd}[1]{\begin{tikzpicture}[baseline=(letter.base)]\node[draw,circle,inner sep=1pt] (letter) {#1};\end{tikzpicture}}
\newcommand{\numeroCerne}[1]{\protect\textcircledd{\bfseries\scriptsize #1}}
\setlist[enumerate,1]{label=\numeroCerne{\arabic*}, leftmargin=*, labelsep=6pt, itemsep=3pt, topsep=2pt, partopsep=0pt, parsep=0pt}

\NewDocumentEnvironment{Exo}{O{breakable} m m m}{%
  \begin{tcolorbox}[
    enhanced,
    #1,
    colback=white,
    colframe=white,
    boxrule=0pt,
    arc=4pt, outer arc=4pt, sharp corners=south,
    left=8pt, right=8pt, top=5pt, bottom=3pt,
    before skip=4pt, after skip=4pt,
  ]
  \noindent\tcbox[on line, colback=skyblue, colframe=skyblue, boxsep=1pt, left=4pt, right=4pt, top=1pt, bottom=1pt, arc=2pt]{\textcolor{white}{\ExoTexte{Exercice~#2}}}~#3\hfill\mbox{\color{skyblue}\normalsize\cursive\hyphenchar\font=-1 #4}\par\vspace{3pt}%
}{%
  \end{tcolorbox}
}

\NewDocumentEnvironment{ExoCont}{O{unbreakable}}{%
  \begin{tcolorbox}[
    enhanced,
    #1,
    colback=white,
    colframe=white,
    boxrule=0pt,
    arc=4pt, outer arc=4pt, sharp corners=south,
    left=8pt, right=8pt, top=5pt, bottom=3pt,
    before skip=4pt, after skip=4pt,
  ]
}{%
  \end{tcolorbox}
}

\newcommand{\Objectif}[2]{%
\vspace{4pt}
\noindent{\large\color{objorange}\bfseries Objectif #1~:}\\{\large\color{objorange}#2}\par
\vspace{2pt}
}

\NewDocumentCommand{\FeuilletTitre}{m}{%
  \begin{center}
  \begin{tcolorbox}[
    enhanced,
    width=0.58\linewidth,
    colback=white,
    colframe=black,
    boxrule=0.8pt,
    arc=8pt,
    sharp corners=downhill,
    left=10pt, right=10pt, top=8pt, bottom=6pt,
    halign=center,
    overlay={
      \node[fill=white, anchor=west, font=\normalsize, inner xsep=4pt]
        at ([xshift=14pt]frame.north west) {Feuille d'exercices n\textdegree{}3};
    },
  ]
  {\huge #1}
  \end{tcolorbox}
  \end{center}
}

\pagestyle{empty}
\begin{document}

\FeuilletTitre{Vecteurs du plan}



\Objectif{5}{Produit d'un vecteur par un nombre r\'eel et colin\'earit\'e de deux vecteurs}

\begin{Exo}[breakable]{1}{\starfull{diffblue}\starfull{diffblue}\starempty}{Calculer}
\renewcommand{\Pt}[1]{\textsf{#1}}
Exprimer le vecteur $\Vj{AC}$ en fonction du vecteur $\Vj{AB}$ dans chacun des cas suivant~:


\noindent\begin{minipage}[t]{0.48\linewidth}\vspace{0pt}
\noindent\numeroCerne{1}\par
{\centering
\begin{tikzpicture}[x=1.88cm,y=0.5cm]
\draw (-2,0) -- (3,0);
\foreach \x in {-1,-0.5,0.5,1,1.5,2.5} \draw (\x,-0.08) -- (\x,0.08);
\node[ptcross] at (0,0) {}; \node[above] at (0,0) {\Pt{A}}; \node[ptcross] at (2,0) {}; \node[above] at (2,0) {\Pt{B}}; \node[ptcross] at (-1.5,0) {}; \node[above] at (-1.5,0) {\Pt{C}};
\end{tikzpicture}\par}
Les vecteurs $\Vj{AC}$ et $\Vj{AB}$ ont m\^eme direction.\\[4pt]
Les vecteurs $\Vj{AC}$ et $\Vj{AB}$ ont~: $\square$~m\^eme sens \quad $\square$~sens oppos\'e.\\[10pt]
On a~: $AC=\dots$ et $AB=\dots$.\\[10pt]
Donc la longueur $AC$ vaut $\dfrac{\dots}{\dots}$ celle de $AB$.\\[10pt]
On a donc~: $\Vj{AC}=\dots\,\dfrac{\dots}{\dots}\,\Vj{AB}$
\end{minipage}\hfill
\begin{minipage}[t]{0.48\linewidth}\vspace{0pt}
\noindent\numeroCerne{2}\par
{\centering
\begin{tikzpicture}[x=2.69cm,y=0.5cm]
\draw (-1,0) -- (2.5,0);
\foreach \x in {-0.5,1,1.5} \draw (\x,-0.08) -- (\x,0.08);
\node[ptcross] at (0,0) {}; \node[below] at (0,0) {\Pt{A}}; \node[ptcross] at (0.51,0) {}; \node[above] at (0.51,0) {\Pt{B}}; \node[ptcross] at (2,0) {}; \node[above] at (2,0) {\Pt{C}};
\end{tikzpicture}\par}
Les vecteurs $\Vj{AC}$ et $\Vj{AB}$ ont m\^eme direction.\\[4pt]
Les vecteurs $\Vj{AC}$ et $\Vj{AB}$ ont~: $\square$~m\^eme sens \quad $\square$~sens oppos\'e.\\[10pt]
On a~: $AC=\dots$ et $AB=\dots$.\\[10pt]
Donc la longueur $AC$ vaut $\dfrac{\dots}{\dots}$ celle de $AB$.\\[10pt]
On a donc~: $\Vj{AC}=\dots\,\dfrac{\dots}{\dots}\,\Vj{AB}$
\end{minipage}

\medskip

\noindent\begin{minipage}[t]{0.48\linewidth}\vspace{0pt}
\noindent\numeroCerne{3}\par
{\centering
\begin{tikzpicture}[x=2.94cm,y=0.5cm]
\draw (-2,0) -- (1.2,0);
\foreach \x in {-1,-0.5,1} \draw (\x,-0.08) -- (\x,0.08);
\node[ptcross] at (0,0) {}; \node[below] at (0,0) {\Pt{A}}; \node[ptcross] at (0.51,0) {}; \node[above] at (0.51,0) {\Pt{B}}; \node[ptcross] at (-1.5,0) {}; \node[above] at (-1.5,0) {\Pt{C}};
\end{tikzpicture}\par}
\end{minipage}\hfill
\begin{minipage}[t]{0.48\linewidth}\vspace{0pt}
\noindent\numeroCerne{4}\par
{\centering
\begin{tikzpicture}[x=3.13cm,y=0.5cm]
\draw (-1,0) -- (2,0);
\foreach \x in {0.5,1.5} \draw (\x,-0.08) -- (\x,0.08);
\node[ptcross] at (0,0) {}; \node[above] at (0,0) {\Pt{A}}; \node[ptcross] at (1,0) {}; \node[above] at (1,0) {\Pt{B}}; \node[ptcross] at (-0.5,0) {}; \node[below] at (-0.5,0) {\Pt{C}};
\end{tikzpicture}\par}
\end{minipage}

\medskip

\noindent\begin{minipage}[t]{0.48\linewidth}\vspace{0pt}
\noindent\numeroCerne{5}\par
{\centering
\begin{tikzpicture}[x=1.71cm,y=0.5cm]
\draw (-1,0) -- (4.5,0);
\foreach \x in {-0.5,0.5,1.5,2,2.5,3,4} \draw (\x,-0.08) -- (\x,0.08);
\node[ptcross] at (0,0) {}; \node[above] at (0,0) {\Pt{A}}; \node[ptcross] at (3.5,0) {}; \node[above] at (3.5,0) {\Pt{B}}; \node[ptcross] at (1,0) {}; \node[above] at (1,0) {\Pt{C}};
\end{tikzpicture}\par}
\end{minipage}\hfill
\begin{minipage}[t]{0.48\linewidth}\vspace{0pt}
\noindent\numeroCerne{6}\par
{\centering
\begin{tikzpicture}[x=1.88cm,y=0.5cm]
\draw (-1,0) -- (4,0);
\foreach \x in {-0.5,1,1.5,2,2.5,3.5} \draw (\x,-0.08) -- (\x,0.08);
\node[ptcross] at (0,0) {}; \node[above] at (0,0) {\Pt{A}}; \node[ptcross] at (3,0) {}; \node[above] at (3,0) {\Pt{B}}; \node[ptcross] at (0.5,0) {}; \node[below] at (0.5,0) {\Pt{C}};
\end{tikzpicture}\par}
\end{minipage}

\end{Exo}

\begin{Exo}[breakable]{2}{\starfull{diffblue}\starfull{diffblue}\starempty}{Repr\'esenter}
\renewcommand{\Pt}[1]{\textsf{#1}}
Dans chacun des cas suivant, repr\'esenter le vecteur $\Vj{AC}$ tel que~:

\noindent\begin{minipage}[t]{0.48\linewidth}\vspace{0pt}
\noindent\numeroCerne{1}~$\Vj{AC}=-\dfrac32\Vj{AB}$\par
{\centering
\begin{tikzpicture}[x=1.88cm,y=0.5cm]
\draw (-2,0) -- (3,0);
\foreach \x in {-1.5,-1,-0.5,0.5,1.5,2,2.5} \draw (\x,-0.08) -- (\x,0.08);
\node[ptcross] at (0,0) {}; \node[above] at (0,0) {\Pt{A}}; \node[ptcross] at (1,0) {}; \node[above] at (1,0) {\Pt{B}};
\end{tikzpicture}\par}
Les vecteurs $\Vj{AC}$ et $\Vj{AB}$ ont m\^eme direction.\\[4pt]
Les vecteurs $\Vj{AC}$ et $\Vj{AB}$ ont~: $\square$~m\^eme sens \quad $\square$~sens oppos\'e.\\[10pt]
On a~: $AB=\dots$, or $AC=\dfrac{\dots}{\dots}\times AB$.\\[10pt]
Donc la longueur $AC$ vaut $\dfrac{\dots}{\dots}\times\dots=\dots$
\end{minipage}\hfill
\begin{minipage}[t]{0.48\linewidth}\vspace{0pt}
\noindent\numeroCerne{2}~$\Vj{AC}=\dfrac78\Vj{AB}$\par
{\centering
\begin{tikzpicture}[x=1.57cm,y=0.5cm]
\draw (-1.5,0) -- (4.5,0);
\foreach \x in {-1,0,0.5,1,1.5,2,2.5,3,4} \draw (\x,-0.08) -- (\x,0.08);
\node[ptcross] at (-0.5,0) {}; \node[above] at (-0.5,0) {\Pt{A}}; \node[ptcross] at (3.5,0) {}; \node[above] at (3.5,0) {\Pt{B}};
\end{tikzpicture}\par}
Les vecteurs $\Vj{AC}$ et $\Vj{AB}$ ont m\^eme direction.\\[4pt]
Les vecteurs $\Vj{AC}$ et $\Vj{AB}$ ont~: $\square$~m\^eme sens \quad $\square$~sens oppos\'e.\\[10pt]
On a~: $AB=\dots$, or $AC=\dfrac{\dots}{\dots}\times AB$.\\[10pt]
Donc la longueur $AC$ vaut $\dfrac{\dots}{\dots}\times\dots=\dots$
\end{minipage}

\medskip

\noindent\begin{minipage}[t]{0.48\linewidth}\vspace{0pt}
\noindent\numeroCerne{3}~$\Vj{AC}=\dfrac14\Vj{AB}$\par
{\centering
\begin{tikzpicture}[x=2.35cm,y=0.5cm]
\draw (-1,0) -- (3,0);
\foreach \x in {-0.5,0.5,1,1.5,2.5} \draw (\x,-0.08) -- (\x,0.08);
\node[ptcross] at (0,0) {}; \node[above] at (0,0) {\Pt{A}}; \node[ptcross] at (2,0) {}; \node[above] at (2,0) {\Pt{B}};
\end{tikzpicture}\par}
\end{minipage}\hfill
\begin{minipage}[t]{0.48\linewidth}\vspace{0pt}
\noindent\numeroCerne{4}~$\Vj{AC}=-\dfrac12\Vj{AB}$\par
{\centering
\begin{tikzpicture}[x=1.57cm,y=0.5cm]
\draw (-2,0) -- (4,0);
\foreach \x in {-1.5,-1,-0.5,0.5,1,1.5,2,2.5,3,3.5} \draw (\x,-0.08) -- (\x,0.08);
\node[ptcross] at (0,0) {}; \node[above] at (0,0) {\Pt{A}}; \node[ptcross] at (3,0) {}; \node[above] at (3,0) {\Pt{B}};
\end{tikzpicture}\par}
\end{minipage}

\medskip

\noindent\begin{minipage}[t]{0.48\linewidth}\vspace{0pt}
\noindent\numeroCerne{5}~$\Vj{AC}=\dfrac23\Vj{AB}$\par
{\centering
\begin{tikzpicture}[x=2.69cm,y=0.5cm]
\draw (-1,0) -- (2.5,0);
\foreach \x in {-0.5,0.5,1,2} \draw (\x,-0.08) -- (\x,0.08);
\node[ptcross] at (0,0) {}; \node[above] at (0,0) {\Pt{A}}; \node[ptcross] at (1.5,0) {}; \node[above] at (1.5,0) {\Pt{B}};
\end{tikzpicture}\par}
\end{minipage}\hfill
\begin{minipage}[t]{0.48\linewidth}\vspace{0pt}
\noindent\numeroCerne{6}~$\Vj{AC}=-\dfrac43\Vj{AB}$\par
{\centering
\begin{tikzpicture}[x=1.18cm,y=0.5cm]
\draw (-2,0) -- (6,0);
\foreach \x in {-1.5,-1,-0.5,0,0.5,1,1.5,2.5,3,3.5,4,4.5,5.5} \draw (\x,-0.08) -- (\x,0.08);
\node[ptcross] at (2,0) {}; \node[above] at (2,0) {\Pt{A}}; \node[ptcross] at (5,0) {}; \node[above] at (5,0) {\Pt{B}};
\end{tikzpicture}\par}
\end{minipage}

\end{Exo}


\begin{Exo}{3}{\starfull{diffblue}\starfull{diffblue}\starempty}{Calculer}
Soient deux points $A$ et $B$. D\'eterminer la valeur des r\'eels $m$, $n$, $p$ et $q$ tels que~:
\begin{multicols}{2}
\begin{enumerate}[label=(\alph*)]
\item $3\Vj{AB}+\dfrac12\Vj{AB}=m\Vj{AB}$
\item $2(3\Vj{BA})=n\Vj{AB}$
\item $2\Vj{AB}-\Vj{BA}=p\Vj{AB}$
\item $3\Vj{BA}-4\Vj{AB}=q\Vj{AB}$
\end{enumerate}
\end{multicols}
\vspace{1pt}
\end{Exo}

\newpage 

\Objectif{6}{D\'ecomposition d'un vecteur en fonction de deux vecteurs non colin\'eaires}

\noindent\begin{minipage}[t]{0.49\linewidth}\vspace{0pt}
\begin{Exo}[unbreakable]{4}{\starfull{diffgreen}\starempty\starempty}{Repr\'esenter}
Les vecteurs $\vec u$ et $\vec v$ sont donn\'es sur le quadrillage ci-dessous. Construire~:
\begin{enumerate}[label=(\alph*)]
\item le point $M$ tel que $\Vj{AM}=\vec u+2\vec v$~;
\item le point $N$ tel que $\Vj{BN}=2\vec u-\vec v$~;
\item le point $R$ tel que $\Vj{CR}=-\vec u+3\vec v$.
\end{enumerate}
\begin{center}
\begin{tikzpicture}[scale=0.7]
\draw[thin,gray!35] (0,0) grid (11,6);
\draw[-{Stealth},thick,blue] (1,5) -- (3,5) node[midway,above]{$\vec u$};
\draw[-{Stealth},thick,red] (5,5) -- (6,6) node[midway,below right]{$\vec v$};
\coordinate (A) at (1,1); \node[ptcross] at (A) {}; \node[below] at (A) {\Pt{A}};
\coordinate (B) at (4,4); \node[ptcross] at (B) {}; \node[above] at (B) {\Pt{B}};
\coordinate (C) at (9,1); \node[ptcross] at (C) {}; \node[below] at (C) {\Pt{C}};
\end{tikzpicture}
\end{center}
\end{Exo}
\end{minipage}\hfill
\begin{minipage}[t]{0.49\linewidth}\vspace{0pt}
\begin{Exo}[unbreakable]{5}{\starfull{diffgreen}\starempty\starempty}{Repr\'esenter}
Sur la figure ci-dessous, exprimer les vecteurs $\vec{w_1}$, $\vec{w_2}$ et $\vec{w_3}$ en fonction de $\vec u$ et $\vec v$.
\begin{center}
\begin{tikzpicture}[scale=0.9]
\draw[thin,gray!35] (-1.5,-1.5) grid (3.5,3.5);
\draw[-{Stealth},thick] (0,0) -- (1,0) node[midway,below]{$\vec u$};
\draw[-{Stealth},thick] (0,0) -- (0,1) node[midway,left]{$\vec v$};
\draw[-{Stealth},thick,blue] (0,0) -- (3,2) node[midway,above]{$\vec{w_1}$};
\draw[-{Stealth},thick,red] (0,0) -- (-1,3) node[midway,above]{$\vec{w_2}$};
\draw[-{Stealth},thick,teal] (0,0) -- (2,-1) node[midway,below right]{$\vec{w_3}$};
\end{tikzpicture}
\end{center}
\end{Exo}
\end{minipage}

\noindent\begin{minipage}[t]{0.49\linewidth}\vspace{0pt}
\begin{Exo}[unbreakable]{6}{\starfull{diffblue}\starfull{diffblue}\starempty}{Repr\'esenter}
La figure ci-dessous est un pavage form\'e de parall\'elogrammes tous identiques. On pose $\vec u=\Vj{AA'}$ et $\vec v=\Vj{AA''}$.\\
Exprimer les vecteurs $\Vj{AB}$, $\Vj{AC}$, $\Vj{BD}$, $\Vj{DC}$ et $\Vj{EB}$ en fonction de $\vec u$ et $\vec v$.
\begin{center}
\begin{tikzpicture}[x={(1.2cm,0cm)},y={(0.5cm,0.95cm)}]
\foreach \i in {0,...,5} \draw[thin,gray!35] (\i,0) -- (\i,4);
\foreach \j in {0,...,4} \draw[thin,gray!35] (0,\j) -- (5,\j);
\draw[-{Stealth},thick,blue] (0,0) -- (1,0) node[midway,below]{$\vec u$};
\draw[-{Stealth},thick,red] (0,0) -- (0,1) node[midway,left]{$\vec v$};
\node[ptcross] at (0,0) {}; \node[below left] at (0,0) {\Pt{A}};
\node[ptcross] at (1,0) {}; \node[below] at (1,0) {\Pt{A'}};
\node[ptcross] at (0,1) {}; \node[above left] at (0,1) {\Pt{A''}};
\node[ptcross] at (3,1) {}; \node[below right] at (3,1) {\Pt{B}};
\node[ptcross] at (1,3) {}; \node[above left] at (1,3) {\Pt{C}};
\node[ptcross] at (4,2) {}; \node[right] at (4,2) {\Pt{D}};
\node[ptcross] at (0,2) {}; \node[left] at (0,2) {\Pt{E}};
\end{tikzpicture}
\end{center}
\end{Exo}
\end{minipage}\hfill
\begin{minipage}[t]{0.49\linewidth}\vspace{0pt}
\begin{Exo}[unbreakable]{7}{\starfull{diffblue}\starfull{diffblue}\starempty}{Raisonner}
$ABCD$ est un parall\'elogramme et $I$ est le milieu de $[CD]$. On pose $\vec u=\Vj{AB}$ et $\vec v=\Vj{AD}$.
Exprimer le vecteur $\Vj{AI}$ en fonction de $\vec u$ et $\vec v$.
\begin{center}
\begin{tikzpicture}[scale=1.3]
\coordinate (A) at (0,0); \coordinate (B) at (3,0); \coordinate (D) at (0.8,2); \coordinate (C) at (3.8,2);
\coordinate (I) at ($(D)!0.5!(C)$);
\draw (A)--(B)--(C)--(D)--cycle; \draw[dashed] (A)--(I);
\node[below] at (A) {\Pt{A}}; \node[below] at (B) {\Pt{B}}; \node[above] at (D) {\Pt{D}}; \node[above] at (C) {\Pt{C}};
\node[above] at (I) {\Pt{I}};
\end{tikzpicture}
\end{center}
\end{Exo}
\end{minipage}

\begin{Exo}[unbreakable]{8}{\starfull{diffred}\starfull{diffred}\starfull{diffred}}{Raisonner}
\begin{minipage}{5.6cm}
\begin{tikzpicture}[scale=1]
\foreach \i/\lab in {0/A,60/B,120/C,180/D,240/E,300/F}
\coordinate (\lab) at (\i:1.4);
\draw (A)--(B)--(C)--(D)--(E)--(F)--cycle;
\draw (A)--(D); \draw (B)--(E); \draw (C)--(F);
\draw[-{Stealth},thick,blue] (0,0) -- (A) node[midway,below]{$\vec u$};
\draw[-{Stealth},thick,red] (0,0) -- (B) node[midway,above left]{$\vec v$};
\node[below right] at (A) {\Pt{A}}; \node[above right] at (B) {\Pt{B}}; \node[above] at (C) {\Pt{C}};
\node[above left] at (D) {\Pt{D}}; \node[below left] at (E) {\Pt{E}}; \node[below] at (F) {\Pt{F}};
\node[below left,inner sep=1pt] at (0,0) {\Pt{O}};
\end{tikzpicture}
\end{minipage}\hfill
\begin{minipage}{\dimexpr\linewidth-5.9cm}
$ABCDEF$ est un hexagone r\'egulier de centre $O$. On pose $\vec u=\Vj{OA}$ et $\vec v=\Vj{OB}$.
\begin{enumerate}
\item Exprimer les vecteurs $\Vj{OD}$, $\Vj{OC}$, $\Vj{OE}$ et $\Vj{OF}$ en fonction de $\vec u$ et $\vec v$.
\item En utilisant la relation de Chasles, exprimer les vecteurs $\Vj{CB}$, $\Vj{AB}$, $\Vj{DF}$ et $\Vj{CD}$ en fonction de $\vec u$ et $\vec v$ uniquement.\\
\textit{Exemple~:} $\Vj{DB}=\Vj{DO}+\Vj{OB}=\vec u+\vec v$.
\end{enumerate}
\end{minipage}
\end{Exo}

\end{document}
